\documentclass[10pt,a4paper]{amsart}
\usepackage{amsmath,amssymb,amsthm,graphicx}
% Retained preamble for the Stacks Project chapter morphisms.tex (Morphisms of Schemes).
% The chapter file pinned-original-source.tex carries only \input{preamble} on its first line, so no
% preamble text lives inside the chapter itself.  The Stacks Project uses one single shared preamble for
% every chapter; the snapshot stages it verbatim in sources/stacks-chow/preamble.tex.  The block below is
% copied mechanically, line for line, from that file's lines 14-18 (the Xy-pic setup) and lines 154-217
% (the native theorem declarations together with the mathematical macros).  Nothing was added, removed or
% redefined.  The wrapper class and the amsmath/amssymb/amsthm/graphicx base packages are supplied by the
% assembler, so the class/\externaldocument part of the shared preamble is not repeated here.
% For commutative diagrams we use Xy-pic
\usepackage[all]{xy}

% We use 2cell for 2-commutative diagrams.
\xyoption{2cell}
% Theorem environments.
%
\theoremstyle{plain}
\newtheorem{theorem}[subsection]{Theorem}
\newtheorem{proposition}[subsection]{Proposition}
\newtheorem{lemma}[subsection]{Lemma}

\theoremstyle{definition}
\newtheorem{definition}[subsection]{Definition}
\newtheorem{example}[subsection]{Example}
\newtheorem{exercise}[subsection]{Exercise}
\newtheorem{situation}[subsection]{Situation}

\theoremstyle{remark}
\newtheorem{remark}[subsection]{Remark}
\newtheorem{remarks}[subsection]{Remarks}

\numberwithin{equation}{subsection}

% Macros
%
\def\lim{\mathop{\mathrm{lim}}\nolimits}
\def\colim{\mathop{\mathrm{colim}}\nolimits}
\def\Spec{\mathop{\mathrm{Spec}}}
\def\Hom{\mathop{\mathrm{Hom}}\nolimits}
\def\Ext{\mathop{\mathrm{Ext}}\nolimits}
\def\SheafHom{\mathop{\mathcal{H}\!\mathit{om}}\nolimits}
\def\SheafExt{\mathop{\mathcal{E}\!\mathit{xt}}\nolimits}
\def\Sch{\mathit{Sch}}
\def\Mor{\mathop{\mathrm{Mor}}\nolimits}
\def\Ob{\mathop{\mathrm{Ob}}\nolimits}
\def\Sh{\mathop{\mathit{Sh}}\nolimits}
\def\NL{\mathop{N\!L}\nolimits}
\def\CH{\mathop{\mathrm{CH}}\nolimits}
\def\proetale{{pro\text{-}\acute{e}tale}}
\def\etale{{\acute{e}tale}}
\def\QCoh{\mathit{QCoh}}
\def\Ker{\mathop{\mathrm{Ker}}}
\def\Im{\mathop{\mathrm{Im}}}
\def\Coker{\mathop{\mathrm{Coker}}}
\def\Coim{\mathop{\mathrm{Coim}}}

% Boxtimes
%
\DeclareMathSymbol{\boxtimes}{\mathbin}{AMSa}{"02}

%
% Macros for moduli stacks/spaces
%
\def\QCohstack{\mathcal{QC}\!\mathit{oh}}
\def\Cohstack{\mathcal{C}\!\mathit{oh}}
\def\Spacesstack{\mathcal{S}\!\mathit{paces}}
\def\Quotfunctor{\mathrm{Quot}}
\def\Hilbfunctor{\mathrm{Hilb}}
\def\Curvesstack{\mathcal{C}\!\mathit{urves}}
\def\Polarizedstack{\mathcal{P}\!\mathit{olarized}}
\def\Complexesstack{\mathcal{C}\!\mathit{omplexes}}
% \Pic is the operator that assigns to X its picard group, usage \Pic(X)
% \Picardstack_{X/B} denotes the Picard stack of X over B
% \Picardfunctor_{X/B} denotes the Picard functor of X over B
\def\Pic{\mathop{\mathrm{Pic}}\nolimits}
\def\Picardstack{\mathcal{P}\!\mathit{ic}}
\def\Picardfunctor{\mathrm{Pic}}
\def\Deformationcategory{\mathcal{D}\!\mathit{ef}}

% --- appended 2026-10-12 (Chief): the project own definitions of the reference/slogan/history
% display environments, copied mechanically from the shared Stacks preamble lines 8-12
% (\usepackage{verbatim} plus the three \newenvironment lines, which the chapter preamble keeps
% verbatim from the project preamble). Nothing added, removed or redefined: these environments
% carry no mathematical content and the project itself typesets them as comments. Without these
% three lines a counted statement whose own text uses them fails with Environment slogan undefined.
\usepackage{verbatim}
\newenvironment{reference}{\comment}{\endcomment}
\newenvironment{slogan}{\comment}{\endcomment}
\newenvironment{history}{\comment}{\endcomment}
\usepackage{cleveref}
\crefname{originallemma}{original lemma}{original lemmas}
\Crefname{originallemma}{Original lemma}{Original lemmas}
\makeatletter
\newlabel{algebra-lemma-dimension-base-fibre-equals-total}{{00ON}{source}{}{sourcepointer-1}{}}
\newlabel{algebra-lemma-dimension-base-fibre-equals-total@cref}{{[originallemma][0][]00ON}{[1][0][]source}{}{}{}}
\makeatother
\crefname{originallemma}{original lemma}{original lemmas}
\Crefname{originallemma}{Original lemma}{Original lemmas}
\makeatletter
\newlabel{algebra-lemma-dimension-base-fibre-total}{{00OM}{source}{}{sourcepointer-2}{}}
\newlabel{algebra-lemma-dimension-base-fibre-total@cref}{{[originallemma][0][]00OM}{[1][0][]source}{}{}{}}
\makeatother
\crefname{originallemma}{original lemma}{original lemmas}
\Crefname{originallemma}{Original lemma}{Original lemmas}
\makeatletter
\newlabel{lemma-dimension-fibre-at-a-point}{{02FX}{source}{}{sourcepointer-3}{}}
\newlabel{lemma-dimension-fibre-at-a-point@cref}{{[originallemma][0][]02FX}{[1][0][]source}{}{}{}}
\makeatother
\title{The dimension of a fibre at a point is at most the sum of the fibre dimensions along the composite, with equality under flatness}
\author{The Stacks Project Authors}
\date{stacks-project:morphisms.tex:local-snapshot-20261012; GNU FDL 1.2 or later}
\begin{document}
\maketitle
\noindent\textit{Editorial note.} All delivered mathematical text is the author's own, copied from the pinned Stacks Project chapter file pinned-original-source.tex; no statement or proof step was rewritten, reordered or removed.\par\medskip
\noindent\textit{Editorial note.} Following the multi-claim convention for this source, source\_url is the source-level URL https://stacks.math.columbia.edu/ and the source label lemma-dimension-fibre-at-a-point-additive is carried in claim\_key and here rather than in a per-claim tag URL.\par\medskip
\noindent\textit{Editorial note.} The preamble is the chapter's shared Stacks Project preamble, retained as sources/stacks-morphisms/retained-preamble.tex: a line-for-line excerpt of the shared Stacks preamble staged with this snapshot. No macro was invented or redefined, and this card adds no preamble text of its own.\par\medskip
\noindent\textit{Editorial note.} Statement span 5395-5407 and proof span 5409-5422 come from the halfA pool card operations/pool-cards/stacks-morphisms-halfA.json and were re-read in the pinned chapter file before this card was written.\par\medskip
\noindent\textit{Editorial note.} The proof refers to the fibre dimension formula of an earlier statement and to two algebra statements; the references are part of the counted span and are handled as navigation pointers.\par\medskip
\noindent\textit{Editorial note.} 3 cross-references occur inside the counted spans; they are listed in external\_references as navigation pointers only.\par\medskip
\noindent\textit{Allowed standard prerequisites (editorial).}\begin{itemize}
\item The description of the dimension of a fibre at a point as the sum of the dimension of the local ring of the fibre and the transcendence degree of the residue field over the residue field of the base point, cited from the same chapter. Cited at proof lines 5411-5412. Reduces the geometric inequality to an inequality between dimensions of local rings of the two fibres.
\item Additivity of the transcendence degree over a tower of field extensions, and the corresponding inequality and equality statements for dimensions of local rings of fibres over a base, cited from the algebra chapter. Used at proof lines 5409-5410 and 5414-5422. Supplies the field degree identity and the two local algebra statements that finish the proof.
\end{itemize}\medskip
\ifdefined\hypertarget\hypertarget{sourcepointer-1}{}\fi
\noindent\textit{Source reference.} Original lemma 00ON: focused-seeds/AGP/sources/stacks-snapshot-20261011/tags.txt, source lines 901--901. Same-chapter reference used in the retained author text at source line(s) 5421, 5699. It is supplied only as a navigation pointer so that the retained text compiles; the counted item delivers the author's own statement and the author's complete proof, and no part of the referenced item is copied in or substituted for the counted argument.\par\medskip
\ifdefined\hypertarget\hypertarget{sourcepointer-2}{}\fi
\noindent\textit{Source reference.} Original lemma 00OM: focused-seeds/AGP/sources/stacks-snapshot-20261011/tags.txt, source lines 900--900. Same-chapter reference used in the retained author text at source line(s) 5419. It is supplied only as a navigation pointer so that the retained text compiles; the counted item delivers the author's own statement and the author's complete proof, and no part of the referenced item is copied in or substituted for the counted argument.\par\medskip
\ifdefined\hypertarget\hypertarget{sourcepointer-3}{}\fi
\noindent\textit{Source reference.} Original lemma 02FX: focused-seeds/AGP/sources/stacks-snapshot-20261011/tags.txt, source lines 3208--3208. Same-chapter reference used in the retained author text at source line(s) 5412, 5659, 5661, 13444. It is supplied only as a navigation pointer so that the retained text compiles; the counted item delivers the author's own statement and the author's complete proof, and no part of the referenced item is copied in or substituted for the counted argument.\par\medskip
\begin{lemma}
\label{lemma-dimension-fibre-at-a-point-additive}
Let $f : X \to Y$ and $g : Y \to S$ be morphisms of schemes.
Let $x \in X$ and set $y = f(x)$, $s = g(y)$.
Assume $f$ and $g$ locally of finite type.
Then
$$
\dim_x(X_s) \leq \dim_x(X_y) + \dim_y(Y_s).
$$
Moreover, equality holds if $\mathcal{O}_{X_s, x}$ is flat
over $\mathcal{O}_{Y_s, y}$, which holds for example if $\mathcal{O}_{X, x}$
is flat over $\mathcal{O}_{Y, y}$.
\end{lemma}

\begin{proof}
Note that $\text{trdeg}_{\kappa(s)}(\kappa(x)) =
\text{trdeg}_{\kappa(y)}(\kappa(x)) + \text{trdeg}_{\kappa(s)}(\kappa(y))$.
Thus by Lemma \ref{lemma-dimension-fibre-at-a-point} the statement
is equivalent to
$$
\dim(\mathcal{O}_{X_s, x})
\leq
\dim(\mathcal{O}_{X_y, x}) + \dim(\mathcal{O}_{Y_s, y}).
$$
For this see Algebra, Lemma \ref{algebra-lemma-dimension-base-fibre-total}.
For the flat case see
Algebra, Lemma \ref{algebra-lemma-dimension-base-fibre-equals-total}.
\end{proof}

\end{document}
